\documentclass[11pt]{amsart}
\usepackage[margin=1.2in]{geometry}
\usepackage{amsmath,amssymb,amsthm}
\usepackage{xcolor}
\definecolor{linkblue}{RGB}{25,60,130}
\usepackage[colorlinks=true,linkcolor=linkblue,citecolor=linkblue,urlcolor=linkblue]{hyperref}

\newtheorem{theorem}{Theorem}[section]
\newtheorem{lemma}[theorem]{Lemma}
\newtheorem{conjecture}[theorem]{Conjecture}
\theoremstyle{remark}
\newtheorem{remark}[theorem]{Remark}
\theoremstyle{definition}
\newtheorem{observation}[theorem]{Observation}

\newcommand{\N}{\mathbb{N}}

\title[Exception sets for $m/n = 1/x+1/y+1/z$]{A verification of Sierpi\'nski's conjecture to $10^{11}$ and complete exception sets for $m/n = 1/x + 1/y + 1/z$, $4 \le m \le 19$}

\author{Jaidev Shah}
\address{}
\email{shahjaidevn@gmail.com}

\subjclass[2020]{11D68 (primary), 11Y50, 11D72 (secondary)}
\keywords{Egyptian fractions, unit fractions, Erd\H{o}s--Straus conjecture, Sierpi\'nski conjecture, Schinzel conjecture}

\begin{document}

\begin{abstract}
Call $n \ge 2$ an \emph{exception for} $m$ if $m/n$ is not a sum of three
unit fractions. Erd\H{o}s and Straus conjectured that $m=4$ has no
exceptions, Sierpi\'nski conjectured the same for $m=5$, and Schinzel
conjectured that every $m \ge 4$ has only finitely many. Pomerance and
Weingartner recently proved that exceptions must exist for every
sufficiently large $m$, with numerical evidence that every $m \ge 20$
has a prime exception in $(m^2, 2m^2)$. We contribute the complementary
small-$m$ data. First, Sierpi\'nski's conjecture holds for all
$2 \le n \le 10^{11}$, extending the classical bound $1{,}057{,}438{,}801$
by a factor of $94.6$. Second, we determine the complete exception sets
for every $6 \le m \le 19$ over all $n \le 10^9$, primes and composites
alike: $243$ exceptions in total, the largest being $83{,}449$ (for
$m=16$), giving empirical values of Schinzel's threshold $n_m$. Third,
we replicate the Pomerance--Weingartner computation for $20 \le m \le 60$
and determine the exact boundary of their window phenomenon: among
$4 \le m \le 60$, the interval $(m^2, 2m^2)$ contains no exception
precisely for $m \in \{4,5,6,7,8,9,11,19\}$. All exceptional values were
certified non-representable by two independently written programs, and
representability certificates were re-verified in exact rational
arithmetic. Code and data accompany the paper.
\end{abstract}

\maketitle

\section{Introduction}

For an integer $m \ge 4$, consider the Diophantine equation
\begin{equation}\label{eq:main}
\frac{m}{n} \;=\; \frac{1}{x} + \frac{1}{y} + \frac{1}{z},
\qquad x, y, z \in \N.
\end{equation}
We say $n \ge 2$ is an \emph{exception for} $m$ if \eqref{eq:main} has no
solution. Erd\H{o}s and Straus \cite{erdos1950} conjectured in 1948 that
$m = 4$ has no exceptions; the conjecture has been verified for
$n \le 10^{17}$ by Salez \cite{salez2014}, with a recent preprint
reporting $10^{18}$ \cite{mihnea2025}. Sierpi\'nski \cite{sierpinski1956}
conjectured the same for $m = 5$, and Schinzel (see \cite{sierpinski1956})
conjectured that for every $m \ge 4$ there is a threshold $n_m$ such that
no $n \ge n_m$ is an exception. Vaughan \cite{vaughan1970} proved that for
fixed $m$ the number of exceptions up to $N$ is
$O\!\left(N \exp(-c(\log N)^{2/3})\right)$, and Elsholtz and Tao
\cite{elsholtztao2013} obtained detailed statistics for the number of
solutions of \eqref{eq:main} when $m = 4$ and $n$ is prime.

The landscape changed recently: Pomerance and Weingartner
\cite{pomeranceweingartner2026} proved that if $n_m$ exists then
$n_m \ge \exp(m^{1/3+o(1)})$ --- exceptions \emph{must} occur for every
sufficiently large $m$ --- and, in explicit form, that for
$m \ge 6.52 \times 10^9$ some prime $p \in (m^2, 2m^2)$ is an exception
for $m$. Their numerical computations support the much stronger assertion
that every $m \ge 20$ has a prime exception in this window.

The map of problem \eqref{eq:main} thus has three regions: a
conjecturally exception-free zone ($m = 4, 5$, and, as our data suggests,
$6$), a zone of nonempty but apparently finite exception sets
($7 \le m \le 19$), and a zone of guaranteed exceptions ($m \ge 20$,
provably for large $m$). This note charts the middle region and pushes
the verification frontier of the first. Our contributions:

\begin{enumerate}
\item Sierpi\'nski's conjecture holds for all $2 \le n \le 10^{11}$
  (Theorem~\ref{thm:sierpinski}). The best bounds we could locate are the
  classical $n \le 1{,}057{,}438{,}801$ (mid-1960s; cited in Graham's
  survey \cite{graham2013}) and, for the residual class
  $n \equiv 1 \pmod 5$, roughly $10^{10}$ in the recent preprint
  \cite{ghermoul2025}.
\item The complete exception sets for $6 \le m \le 19$ over all
  $n \le 10^9$ (Theorem~\ref{thm:sets}), yielding empirical values of
  $n_m$ that we could not locate in the literature or in the OEIS.
\item An independent replication of the Pomerance--Weingartner window
  computation for $20 \le m \le 60$, together with the exact boundary of
  the phenomenon: among $4 \le m \le 60$ the window $(m^2, 2m^2)$ is
  exception-free precisely for $m \in \{4,5,6,7,8,9,11,19\}$
  (Section~\ref{sec:window}).
\end{enumerate}

\begin{remark}\label{rem:priority}
We were unable to consult the numerical section of
\cite{pomeranceweingartner2026} beyond its abstract while preparing this
note. Should it or its references already contain exception sets for
$m \le 19$, priority for the overlap of course belongs there; the
$10^9$-exhaustiveness (including composite $n$), the $n_m$ table, and the
Sierpi\'nski bound appear in any case to be new.
\end{remark}

\section{The decision criterion}\label{sec:criterion}

All computations rest on the following elementary and well-known
reduction, which we state with proof to make the note self-contained.
The essential point is that it is a \emph{decision procedure}: it can
certify non-representability, not merely fail to find solutions.

\begin{lemma}\label{lem:criterion}
Let $m \ge 4$ and $n \ge 2$ with $m \nmid n$. Then \eqref{eq:main} is
solvable if and only if there exist an integer $x$ with
$n/m < x \le 3n/m$ and a divisor pair $d \cdot d' = b^2$, where
$a = mx - n$ and $b = nx$, such that
\[
d \equiv d' \equiv -b \pmod{a}.
\]
In that case $y = (b+d)/a$ and $z = (b+d')/a$ solve \eqref{eq:main}
together with $x$. If $m \mid n$, then \eqref{eq:main} is solvable, e.g.\
$m/n = 1/(2q) + 1/(3q) + 1/(6q)$ with $q = n/m$.
\end{lemma}

\begin{proof}
Suppose \eqref{eq:main} holds with $x \le y \le z$. Then
$1/x < m/n \le 3/x$, so $n/m < x \le 3n/m$, and subtracting $1/x$ gives
$1/y + 1/z = a/b$ with $a = mx - n \ge 1$, $b = nx$. From
$ayz = b(y+z)$ we get $(ay - b)(az - b) = b^2$; put $d = ay - b$ and
$d' = az - b$, both positive and congruent to $-b$ modulo $a$, with
$dd' = b^2$.

Conversely, given such a pair, $y = (b+d)/a$ and $z = (b+d')/a$ are
positive integers, and
\[
\frac{1}{y} + \frac{1}{z}
= \frac{a\,(2b + d + d')}{(b+d)(b+d')}
= \frac{a\,(2b + d + d')}{b\,(2b + d + d')}
= \frac{a}{b},
\]
using $(b+d)(b+d') = b^2 + b(d+d') + dd' = b(2b + d + d')$. Adding $1/x$
recovers \eqref{eq:main}.
\end{proof}

Enumerating all $x$ in the stated range and, for each $x$, all divisors
$d \le b$ of $b^2$ (obtained from the full factorization of $b = nx$)
therefore decides \eqref{eq:main} for any given pair $(m, n)$.

Two standard reductions structure the search over $n$. First, if
\eqref{eq:main} is solvable for $n$ then it is solvable for every
multiple of $n$ (scale the denominators); hence a composite exception
must have \emph{all} of its divisors $> 1$ exceptional, and in particular
all of its prime factors. Second, $n < m/3$ is trivially an exception
since then $m/n > 3$; we therefore only report exceptions with
$n \ge m/3$, in line with the convention of \cite{pomeranceweingartner2026}.

\section{Computations}\label{sec:computations}

\subsection{Production engine}
A C program (\texttt{mfrac.c} in the accompanying repository) combines a
segmented sieve of Eratosthenes with a first-hit certificate search for
each prime $p$: candidate least denominators
$x = \lfloor p/m \rfloor + 1, \lfloor p/m \rfloor + 2, \dots$ are tried
in order, and for each $x$ the divisor condition of
Lemma~\ref{lem:criterion} is tested over divisors $d = p^{j} f$
($j \in \{0,1,2\}$, $f \mid x^2$) built from a partial factorization of
$x$, with all congruence arithmetic exact in 64/128-bit integers. A prime
not certified within $96$ candidate values of $x$ escalates to $20{,}000$
candidates; anything still uncertified is flagged as a suspect and
\emph{never} silently classified. Suspects are then settled by the full
decision procedure of Section~\ref{sec:criterion}. Across all runs
reported here ($4.6$ billion primes) every suspect was resolved; for
$m = 12$, for instance, the eleven suspects raised by the scan are
precisely the eleven exceptions $433, \dots, 12241$ of
Theorem~\ref{thm:sets}.

For every $m$ the range $n \le 10^6$ (primes \emph{and} composites) was
settled by the complete decision procedure directly. The composite gap
above $10^6$ was closed by the divisor reduction: all
$4{,}769{,}723$ products of exceptional primes lying in $(10^6, 10^9]$
across $7 \le m \le 19$ were decided individually; none is an exception.
For $m = 5$ the composite case is immediate: every composite
$n \le 10^{11}$ has a prime factor below $\sqrt{10^{11}} < 10^6$, and no
$n \le 10^6$ is an exception for $m=5$.

\subsection{Independent verification}
Because a claimed exception is a negative statement, we re-verified every
one with a second program (\texttt{check2.c}) written independently around
a different algorithm: for each admissible $x$ it iterates $y$ over the
full interval $(b/a,\, 2b/a]$ and tests integrality of
$z = by/(ay - b)$, with no divisor enumeration, no factorization, and no
modular shortcuts. All $243$ claimed exceptions were re-certified, with
representable neighbours as positive controls. In addition, a brute-force
checker over exact rationals (Python \texttt{fractions}) was run against
both C programs for all $m \le 14$, $n < 400$ (perfect agreement), and
$200$ randomly sampled prime certificates with $m \in [5,19]$,
$p \le 10^9$ were re-verified in exact rational arithmetic ($200/200$
valid). As a further consistency check, the number of primes processed in
the Sierpi\'nski run, $4{,}117{,}976{,}315$, equals
$\pi(10^{11}) - \pi(10^6)$ exactly.

The total computational cost of all results in this note is under one
CPU-hour on four cores of a commodity cloud container.

\section{Results}\label{sec:results}

\begin{theorem}\label{thm:sierpinski}
For every integer $2 \le n \le 10^{11}$ there exist positive integers
$x, y, z$ with $5/n = 1/x + 1/y + 1/z$.
\end{theorem}

\begin{theorem}\label{thm:sets}
For $6 \le m \le 19$, the exceptions $n \le 10^9$ (with $n \ge m/3$) of
equation \eqref{eq:main} are exactly the following; there are none for
$m = 6$.
\begin{center}
\small
\begin{tabular}{r p{0.78\textwidth}}
\hline
$m$ & exceptions $n \le 10^9$ \\
\hline
$6$ & none \\
$7$ & 2 \\
$8$ & 2, 3, 11, 17, 131, 241 \\
$9$ & 2, 5, 11, 19 \\
$10$ & 2, 3, 7, 11, 43, 61, 67, 181 \\
$11$ & 2, 3, 4, 37 \\
$12$ & 2, 3, 5, 7, 13, 25, 29, 31, 37, 73, 97, 193, 433, 577, 1129, 1657, 1873, 2521, 2593, 3433, 10369, 12049, 12241 \\
$13$ & 2, 3, 4, 5, 7, 14, 53, 61, 67, 79, 211, 281 \\
$14$ & 2, 3, 4, 5, 17, 19, 29, 59, 257, 353, 841 \\
$15$ & 2, 3, 4, 8, 16, 17, 19, 23, 31, 34, 47, 53, 61, 79, 113, 122, 137, 151, 197, 226, 233, 271, 541, 1103, 1171, 1367, 4201, 6301, 12601, 16831, 20521 \\
$16$ & 2, 3, 4, 5, 6, 7, 9, 11, 13, 17, 22, 23, 33, 34, 37, 73, 97, 113, 121, 131, 167, 193, 241, 257, 262, 421, 482, 577, 593, 641, 769, 1201, 1489, 2113, 2521, 2689, 3169, 3361, 4801, 4993, 5281, 8161, 8641, 33601, 36529, 78721, 83449 \\
$17$ & 2, 3, 4, 5, 6, 7, 9, 13, 19, 23, 41, 43, 53, 71, 73, 157, 281, 421, 1123, 2081 \\
$18$ & 2, 3, 4, 5, 7, 10, 11, 13, 19, 22, 23, 29, 31, 37, 38, 41, 47, 59, 61, 73, 109, 113, 131, 137, 149, 181, 193, 223, 239, 281, 379, 389, 397, 433, 457, 541, 599, 613, 661, 761, 811, 821, 911, 1009, 1297, 1381, 2269, 2819, 9461, 16561, 17389, 28081, 35281 \\
$19$ & 2, 3, 4, 5, 6, 7, 8, 10, 11, 13, 23, 26, 29, 41, 43, 46, 82, 97, 137, 181, 193, 229, 353 \\
\hline
\end{tabular}
\end{center}
\end{theorem}

Under Schinzel's conjecture that each exception set is finite, and on the
evidence that each set above terminates more than four orders of
magnitude below the search bound, the resulting empirical thresholds are
\[
\begin{array}{c|cccccccccccccc}
m & 6 & 7 & 8 & 9 & 10 & 11 & 12 & 13 & 14 & 15 & 16 & 17 & 18 & 19\\
\hline
n_m & 2 & 3 & 242 & 20 & 182 & 38 & 12242 & 282 & 842 & 20522 & 83450 & 2082 & 35282 & 354
\end{array}
\]
These values are exact if and only if no exception exceeds $10^9$;
Pomerance--Weingartner's lower bound $n_m \ge \exp(m^{1/3+o(1)})$
concerns large $m$ and does not by itself preclude further exceptions
for $m \le 19$, so the values above should be read as conjectural.

\begin{observation}\label{obs:structure}
(i) Composite exceptions occur only for $m \ge 11$, and in exactly three
cases a prime and its square are both exceptional for the same $m$:
$(5, 25)$ for $m = 12$, $(29, 841)$ for $m = 14$, and $(11, 121)$ for
$m = 16$.

(ii) The largest exceptions concentrate in the residue class
$n \equiv 1 \pmod m$. All eleven exceptions above $400$ for $m = 12$ are
$\equiv 1 \pmod{24}$ --- the classes in which the Erd\H{o}s--Straus
equation resists all polynomial identities --- and the four largest
exceptions for $m = 18$, as well as the three largest for $m = 15$, are
all $\equiv 1 \pmod m$. The largest exception overall, $83449$ for
$m = 16$, satisfies $83449 \equiv 9 \pmod{16}$, again a square class.
This is consistent with the quadratic-residue obstruction of Schinzel in
the $m = 4$ theory (see \cite{mordell1969, elsholtztao2013}).
\end{observation}

\section{The window boundary}\label{sec:window}

Running the complete decision procedure over all primes
$p \in (m^2, 2m^2)$ for $20 \le m \le 60$ reproduces the numerical claim
of \cite{pomeranceweingartner2026}: every such $m$ has at least one prime
exception in its window, $1494$ in total. (The counts grow roughly
linearly, from $6$ at $m = 20$ to $120$ at $m = 60$.) Since
Theorem~\ref{thm:sets} determines the small-$m$ exception sets
completely, the boundary of the phenomenon can be stated exactly.

\begin{theorem}\label{thm:window}
Among $4 \le m \le 60$, the interval $(m^2, 2m^2)$ contains no exception
for $m$ precisely when
\[
m \in \{4,\ 5,\ 6,\ 7,\ 8,\ 9,\ 11,\ 19\}.
\]
\end{theorem}

In particular the window phenomenon described in
\cite{pomeranceweingartner2026} for $m \ge 20$ in fact begins at
$m = 10$, with exactly two later gaps at $m = 11$ and $m = 19$. The data
suggests:

\begin{conjecture}\label{conj:window}
For every $m \ge 4$ outside $\{4,5,6,7,8,9,11,19\}$, the interval
$(m^2, 2m^2)$ contains a prime exception for $m$.
\end{conjecture}

For $m \ge 20$ this is supported by \cite{pomeranceweingartner2026} (and
proved there for $m \ge 6.52\times 10^9$); the content added by our data
is that the eight listed values are the \emph{only} failures.

\section{Data and code availability}

The two C programs (\texttt{mfrac.c}, \texttt{check2.c}), the complete
exception lists, the window-exception lists for $20 \le m \le 60$, and a
run-by-run verification summary are available at
\begin{center}
\url{https://github.com/shahjaidev/explainers/tree/main/unit-fractions}
\end{center}
Every claimed representation can be re-verified from the published
witnesses in exact arithmetic, and every claimed exception is
reproducible with either program in seconds.

\section*{Acknowledgments}

The computations, code, and a draft of this note were produced with the
assistance of an AI system (Claude, Anthropic), working under the
author's direction; the author reviewed and takes responsibility for all
claims. We thank the authors of \cite{pomeranceweingartner2026}, whose
theorem motivated this computation.

\begin{thebibliography}{10}

\bibitem{elsholtztao2013}
C.~Elsholtz and T.~Tao,
\emph{Counting the number of solutions to the Erd\H{o}s--Straus equation
on unit fractions},
J. Aust. Math. Soc. \textbf{94} (2013), 50--105.

\bibitem{erdos1950}
P.~Erd\H{o}s,
\emph{Az $1/x_1 + 1/x_2 + \dots + 1/x_n = a/b$ egyenlet eg\'esz sz\'am\'u
megold\'asair\'ol},
Mat. Lapok \textbf{1} (1950), 192--210.

\bibitem{ghermoul2025}
B.~Ghermoul,
\emph{Almost a complete proof of the generalized Erd\H{o}s--Straus
conjecture: $5/a = 1/b + 1/c + 1/d$},
arXiv:2508.07367 (2025).

\bibitem{graham2013}
R.~L. Graham,
\emph{Paul Erd\H{o}s and Egyptian fractions},
in: Erd\H{o}s Centennial, Bolyai Soc. Math. Stud. \textbf{25},
Springer, 2013, 289--309.

\bibitem{mihnea2025}
S.~Mihnea and D.~C. Bogdan,
\emph{Further verification and empirical evidence for the
Erd\H{o}s--Straus conjecture},
arXiv:2509.00128 (2025).

\bibitem{mordell1969}
L.~J. Mordell,
\emph{Diophantine Equations},
Academic Press, London, 1969.

\bibitem{pomeranceweingartner2026}
C.~Pomerance and A.~Weingartner,
\emph{Exceptions to the Erd\H{o}s--Straus--Schinzel conjecture},
Ramanujan J. \textbf{69} (2026); arXiv:2511.16817.

\bibitem{salez2014}
S.~E. Salez,
\emph{The Erd\H{o}s--Straus conjecture: new modular equations and
checking up to $N = 10^{17}$},
arXiv:1406.6307 (2014).

\bibitem{sierpinski1956}
W.~Sierpi\'nski,
\emph{Sur les d\'ecompositions de nombres rationnels en fractions
primaires},
Mathesis \textbf{65} (1956), 16--32.

\bibitem{vaughan1970}
R.~C. Vaughan,
\emph{On a problem of Erd\H{o}s, Straus and Schinzel},
Mathematika \textbf{17} (1970), 193--198.

\end{thebibliography}

\end{document}
